\section{Uniform definitions and common closure levels}
\label{close:section-uniform}

The arguments above use two different finite inductions.  The first
iterates marked increments, with a fixed prescription in the root sort.
The second ranges over finite geometric shapes.  In each application we
need one cardinal at which every finite stage is available.  We now give
the definability and closure bookkeeping that permits this in ZFC.
The fixed point data, and subsequently the threshold furnished by the
marked-increment application, are fixed before performing the bookkeeping
for the geometric application.  Thus the two applications may use different
closure functions.

The bookkeeping has three stages.  First we express the finite systems
and their required properties by fixed formulas at a given cardinal.
Next we list the bound requests and the order in which each construction
uses them.  Only after those formulas are fixed do we choose a common
majorant and construct the nested eligible subclasses introduced in the
\hyperref[lev:closure-preview]{preview in Section~\ref*{lev:section-working-schemes}}.
Their absorption property will allow
each finite induction to retain an unbounded supply of levels.

\subsection{Set definitions at a given cardinal}

\begin{proposition}[Uniform finite-depth definitions]
\label{close:uniform-definitions}
Fix the set code for $\K$ from
\Cref{found:uniform-reduction}, its order template, and the fixed
data used in either of the two applications.  There are a set $\mathcal D$
of normalized finite descriptors and fixed first-order formulas of set
theory with the following properties.
\begin{enumerate}
\item Given $d\in\mathcal D$ and an infinite cardinal $\mu$, the formulas
define sets of standardized objects and arrows.  For marked descriptors
these sets are obtained by a finite recursion whose length is an input;
at eligible levels they give precisely the categories $\mathcal C_n$.
For a geometric descriptor they give the systems $\mathcal G_L$ and their
arrows.
\item The formulas define orbit comparison, comparison of realized types
over a base, relative isolation, the prescribed root marginal, and the
finite projection and boundary conditions.  On levels with the requisite
working package these definitions agree with their mathematical meanings
in the preceding sections.
\item There is one unary adequacy formula, with a descriptor as parameter,
which can assert any of the packages needed during the finite inductions:
the individual working conditions, the chart clause using $D_F$, the
projection and type conditions, amalgamated completion, exact boundary
lifting, or the geometric categoricity and basis conditions.  Its use
does not require a formula for a previously chosen closure function.
\item All candidate term templates, named patterns, and finite descriptors
needed in the two inductions can be chosen from $\mathcal D$.  In
particular, their codes have one set-sized bound independent of the
finite depth and of the level cardinal.
\end{enumerate}
These are definitions at every cardinal; a definition by itself makes no
claim that the cardinal is eligible.
\end{proposition}

\begin{proof}
We specify the set recursion and the formulas used in it.  The small-model
code for $\K$ defines membership and strong substructure by the directed
union formulas in \Cref{found:uniform-reduction}.  These are fixed
formulas with a set parameter.  They are therefore available in all the
following set constructions.

\paragraph{Standardized structures and tuples.}
For a fixed finite descriptor and cardinal $\mu$, use a separate copy of
$\mu$ for each vertex sort, and represent each inclusion by its map.
Pointed data and marks are represented by constants in their indicated
sorts.  This convention allows a proper same-cardinality inclusion: it is
an injective map between the two copies, rather than the identity map on
a common universe.  A finite diagram of such structures is specified by
sets of finitary relations, function graphs, constant values, and maps
between these carriers.  The collection of all candidates is a set.
An arbitrary diagram of the specified cardinalities has a copy in this
collection, by choosing bijections in its finitely many vertex sorts.

The bounded tuples used below are functions with ordinal domain of
cardinality less than $\mu$, together with a vertex occurrence for every
coordinate.  It suffices to use domains that are cardinals below $\mu$:
other enumerations are obtained by reindexing.  All such functions form
a set.  A fragment in a system is likewise an occurrence-tagged subset
of its finitely many carriers.  Equalities between occurrences are tested
after their specified maps into a common vertex; they are not inferred
from equality of sort tags.

\paragraph{Raw comparison relations.}
Suppose sets $\mathcal O$ and $\mathcal A$ of objects and arrows have already
been constructed.  Isomorphism is expressed by an arrow with a two-sided
inverse in $\mathcal A$.  Put
\[
 \operatorname{Orb}_{\mathcal O,\mathcal A}(Y,u;Z,v)
 \quad\Longleftrightarrow\quad
 \exists h\colon Y\mathrel{\cong}Z\ (h(u)=v).
\]
Here, and in the formulas that follow, the quantified objects and arrows
belong to the displayed sets.  If $e\colon X\to Y$ and $f\colon X\to Z$
are arrows, put
\begin{align*}
 &\operatorname{Eq}_{\mathcal O,\mathcal A}
       (X;e,Y,u;f,Z,v)\\
 &\qquad\Longleftrightarrow
 \exists W\ \exists g\colon Y\to W\ \exists h\colon Z\to W
       \bigl(g e=h f\ \text{ and }\ g(u)=h(v)\bigr).
\end{align*}
These are formulas about sets.  At a level where amalgamation has been
proved, the second relation is precisely equality of the realized types
used in the one-dimensional calculus, as characterized in
\Cref{calc:scheme-amalgamation}.  Until amalgamation has been
established we use the displayed relation itself, without taking a
quotient or claiming transitivity.

Here is the isolation test in full.  For $e\colon X\to Y$ and bounded
$b,c$ in $Y$, the test is
\begin{align*}
 \exists z\in X^{<\mu}\ \forall (f\colon X\to Z)\ \forall b',c'\quad
 \bigl[&\operatorname{Eq}(X;e,Y,b;f,Z,b')\\[-2pt]
 &{}\land\operatorname{Orb}
       (Y,(b,c,e z);Z,(b',c',f z))\bigr]\\
 &\hspace{12mm}\Longrightarrow
       \operatorname{Eq}(X;e,Y,(b,c);f,Z,(b',c')).
\end{align*}
The tuples $b',c'$ have the same lengths and occurrences as $b,c$.
The universal quantifiers range over all extensions in $\mathcal O$,
not just over the objects about to be defined.  The displayed notation
$X^{<\mu}$ includes the occurrence tags.  By the diagram and amalgamation
clauses of the working package, this is exactly isolation relative to
$b$: $z$ enumerates an isolating fragment.  Requiring this formula for
every bounded $c$ in $Y$ defines all-domination over $X+b$.  In particular,
this test uses only $\mathcal O$ and $\mathcal A$ of the preceding
category.

\paragraph{Root prescriptions and free lifting.}
The base descriptor includes the fixed named pattern.  A prescription
for a root type is given by the approximation code already used in
Lemma~\ref{geom:root-prescription}: a named bounded list of approximating
tuples, its joint underlying $\K$-diagram with the point, and a test
arity sufficient to determine the type.  For any test at that arity,
the requested comparison is the comparison given by all but the allowed
number of exceptions in the list.  The allowed number is
$|z|+\aleph_0$ for a parameter test $z$, and the list is longer than this
number.  Checking this rule for every such $z$ is a fixed formula using
the underlying $\K$-diagram predicate.  Existence and uniqueness of a
type satisfying it, and its free lifting, are the conclusions proved
when the prescription is introduced; they are not part of the assertion
that the formula is defined.

This root test is applied after forgetting to the root $\K$-sort.
The projection type correspondence in
\Cref{iter:projection-types} identifies it, at the levels under
consideration, with the required type over the whole system.  Thus the
definition does not successively quote a new formula for a type at each
marked depth.

For completeness, references to free lifting in an adequacy condition
also have a fixed set interpretation.  Let $p$ be represented by
$(A\to Y,u)$, and let $q$ be represented by $(B\to Z,v)$ over an extension
$j\colon A\to B$, with $q$ restricting to $p$.  This is a same-level
test: $A,B,Y,Z$ are objects at $\mu$.  Mixed lifting is derived below
using approximation costs bounded at the source level.  For an infinite
$\sigma\ge |u|$, an approximation test on
$I=(u_i)_{i<\theta}$ requires the following.  For each parameter tuple
$z$ of length at most $\sigma$ in the base, and each corresponding
subtuple of $u$ and the $u_i$, at most $|z|+\aleph_0$ indices have a
joint diagram with $z$ different from that of the realized type.
The free-lift test is
\[
 \forall\sigma\ (|u|+\aleph_0\le\sigma<\mu)\quad
 \exists\theta\ (\sigma<\theta<\mu)\quad
 \exists I\in(A^{|u|})^\theta
\]
such that $I$ passes this test for $p$ over $A$ and $jI$ passes it
for $q$ over $B$.  Here $\sigma,\theta$ are cardinals, and all joint
diagrams are computed in the indicated realizing objects.
The approximation construction gives these sequences for a free lift.
Conversely, compare a sequence witnessing this test with the sequence
computing the free lift at the same arity.  Move a parameter query from
$B$ into $A$ over both sequences by bounded saturation.  Each majority
then computes the original type on the moved query; the two majorities
agree because their sequences have lengths exceeding their exception
bounds.  This proves equality with the free lift on every bounded test,
and hence equality of types by \Cref{calc:scheme-amalgamation}.
All the sequences, tuples, exception sets, and realized types in this
test range over sets at the given cardinal.  This supplies a formula
without a satisfaction predicate for a category definition.

\paragraph{The marked recursion.}
Let $\mathcal O_0(\mu)$ be the standardized pointed $\K$-models with the
specified named pattern, and let $\mathcal A_0(\mu)$ be their strong maps
preserving the names.  Given $\mathcal O_i(\mu),\mathcal A_i(\mu)$, form
their raw comparison and isolation relations as above, and set
\[
\begin{split}
\mathcal O_{i+1}(\mu)=
 \{(e\colon X\to Y,b):\;&X,Y\in\mathcal O_i(\mu),\quad
                         e\in\mathcal A_i(\mu),\\
 &b\text{ passes the prescribed root test},\\
 &\forall c\in Y^{<\mu}\quad
               (b,c)/X\text{ passes the relative isolation test}\}.
\end{split}
\]
An arrow from $(e\colon X\to Y,b)$ to
$(e'\colon X'\to Y',b')$ is a pair
$(f\colon X\to X',g\colon Y\to Y')$ of arrows in $\mathcal A_i(\mu)$
with $g e=e'f$ and $g(b)=b'$.  This defines
$\mathcal A_{i+1}(\mu)$ by Separation.  Reading a pair as two layers
gives the next axis; the earlier axes, their maps, their marks, and the
root projection are part of its finite diagram code.

The operator in these two definitions is fixed.  A finite sequence
\[
 (\mathcal O_i(\mu),\mathcal A_i(\mu))_{i\le n}
\]
satisfying the initial condition and this operator is a set and is
unique.  Existence follows by finite recursion using Power Set and
Separation at each step.  Consequently the assertion that a candidate
belongs to its last object or arrow set is one first-order formula with
$n$ as an input.  It is not a sequence of truth definitions on the
universe.  Induction, using the prescribed-type correspondence and
\Cref{mark:package}, identifies this recursion with the intended
$\mathcal C_n$ at eligible levels.  Raw tests at other levels have no role
in that identification.

\paragraph{Geometric descriptors.}
Such a descriptor contains a finite axis set and its active-label family.
The shape condition is a finite calculation.  Membership in $P$, and
independence of a finite list in $P$, are tested in a strong
$\lambda_0$-section over $M_0$, as in
\Cref{geom:pregeometry}.  A section, an isomorphism of sections,
and a finite product witness are set-coded data.  The independence
predicate so obtained is fixed, with $M_0,p_0,\lambda_0$ as parameters.
For a set $H$ in a vertex, independence means that every finite subset
is independent; spanning means that every relevant point is in the
closure of some finite subset of $H$.  Thus the existence of pieces
$(H_s)_{s\in L}$ forming the required partitioned basis is a first-order
condition on subsets of the finitely many carriers.  Their required
cardinalities are expressed by bijections.

The arrow condition is componentwise strongness together with
preservation of finite independence in a vertex modulo its proper lower
vertices.  A dependence modulo the lower vertices has a finite witness
there.  This gives a fixed formula, equivalent to extendibility of basis
pieces by \Cref{sys:shapes}.  The bases are existential witnesses
in these tests; they are not extra structure that an isomorphism must
preserve.  Inactive vertices and the equalities forced by their largest
active lower label are included as part of the system condition.  This
constructs the geometric object and arrow sets without a recursion in
formula complexity.

\paragraph{A set of template codes.}
Choose an infinite cardinal $b_0$ bounding the fixed vocabulary, template
labels, point code, and the lengths and pattern codes allowed for the
named root prescriptions.  The latter have such a bound because the
approximation list length is bounded in terms of the underlying
$\K$-testing arity, independently of the level where an actual pair is
found.  Order and equality patterns for lists of this fixed bounded
length form a set.

At one marked-template extraction, the expansion consists of old chart
labels and representation tags, finitely indexed coordinate functions,
the finite system maps and mark information, a retraction-index sort
with its evaluation map, and the vertex LS-witness functions.  For a
vocabulary and label bound $b$, these symbols have a cardinal bound
depending only on $b$ and the fixed data.  There is no symbol for every
element of a sample or every entry of a sample-sized diagram table:
sample entries and retraction indices are elements of sorts.  Complete
atomic diagrams on finite lists of terms are subsets of the set of
atomic term instances and their negations.

For definiteness, increase bounds generously and take
\[
 b_{m+1}=\bigl(2^{\,2^{b_m}}\bigr)^+,
 \qquad b_* =\sup_{m<\omega}b_m.
\]
Induction bounds the languages, finite-term sets, and their possible
diagram codes at each finite extraction by a corresponding $b_m$.
Normalize symbol names by ordinals below these bounds.  The union over
$m<\omega$ of the resulting sets of candidate codes is a set containing
every candidate template needed at finite depth.  The explicit geometric
tile templates obey the same bound.  Finally take finite diagrams of
these codes, together with finite axis data, root codes, and the flags
described next.  This is the promised set $\mathcal D$.

\paragraph{One adequacy formula.}
Use the following fixed menu of assertions.  A descriptor can request
any finite selection of them for any of its indicated lower schemes.
Every selection, including a preliminary geometric categoricity or
basis selection, requires a strong limit cardinal of countable
cofinality above all its fixed parameter bounds.
\begin{enumerate}
\item For every well-order
$I$ of cardinality $\mu$, the indicated term interpretation $F(I)$ is
an object, every object at the level is isomorphic to $F(I)$, and
substitution along every same-level order embedding gives an arrow.
The term interpretation is its set of terms modulo the template
equality relation.  It suffices to quantify over orders on a fixed
copy of $\mu$, so this is a first-order test on sets.  In particular,
validity of only the presentation on the initial ordinal $\mu$ is
not the presentation clause used here.
\item Same-level arrows preserve bounded tuple orbits, and comparison of
bounded tuples in well-ordered charts is exactly the fixed comparison
$D_F$ of \Cref{lev:diagram-coding}.
\item A chain of objects whose union has the level cardinality has an
object as its union.  For this test it suffices to use a chain with at
most $\mu$ indices: from any chain with a union of size $\mu$, choose
stages covering its elements in all vertex sorts.  The resulting
subchain has the same union and has order type below $\mu^+$.
\item The root prescription, the indicated projections on objects and
arrows, diagram comparison on projection sorts, and the stated type
correspondence over a base hold.  The comparisons use the raw relations
above, now together with the requested working conditions.  The
free-lift clause uses the approximation test just specified.
\item A compatible boundary arrow has an amalgamated completion, or a
compatible boundary isomorphism extends to an isomorphism of the whole
objects, as indicated by the flag.  These assertions quantify over the
finite families of maps on the proper axis subsets and the object and
arrow sets at this cardinal.
\item For a geometric descriptor, the indicated categoricity, lifting
with prescribed nonroot basis bijections and root map, or lifting over
a fixed base with prescribed complementary basis bijections holds.
The pieces and the bijections range over sets of subsets and functions
of the standardized carriers.
\item A designated term interpretation $T$ into an indicated lower
scheme has the following properties.  Its values on the orders used
by the specified natural layout are objects of that lower scheme,
its same-level order substitutions give lower-scheme arrows, and
comparison of its bounded term lists agrees with the lower scheme's
$D_F$.  Require explicitly that this comparison is natural in the
new interpretation: two instances with the same $T$-labels and joint
horizontal support pattern have the same lower-scheme diagram, and
the descriptor's support substitutions preserve the designated term
parts.  This is the comparison used for the natural-layout test;
arbitrary old-chart representations without this horizontal
naturality assertion do not suffice.  Orders, their finitely many
designated parts, term evaluations, and support patterns are set-coded
at the level.  This flag may be
requested for the base and upper interpretations of a new marked
pair template while omitting the new pair's chain and chart flags.
It uses the already established lower-scheme comparison.
\end{enumerate}
There are only finitely many forms in this menu.  The geometric
categoricity and basis item can be used before the diagram clause has
been obtained for the current geometric shape.
A descriptor records which assertions are currently required, and for
which finitely many lower descriptors.  Bounded quantification over that
finite record yields one formula, denoted
$\operatorname{Adeq}(d,\mu)$.  Its ingredients were all defined by set
operations or by the fixed underlying $\K$ and $D_F$ formulas.  This
proves all four assertions.
\end{proof}

The definitions above describe objects and assertions at one cardinal;
they do not themselves establish eligibility there.  Before collecting
the bounds, we check how the individual adequacy conditions supply the
maps and comparisons between different eligible levels that the earlier
constructions use.

\begin{lemma}[The mixed-level interface]
\label{close:mixed-interface}
For either application, suppose the relevant individual adequacy
conditions of Proposition~\ref{close:uniform-definitions} hold on a class
$S$ of levels, above the fixed parameter bounds.  When the full individual
working conditions, including the chart clause, are requested, the mixed
presentation, coherence, LS, diagram, saturation, and map conditions
needed in the preceding arguments hold on $S$.  A preliminary descriptor
requests only the corresponding already established clauses: in
particular, its natural-interpretation flag supplies comparisons in the
indicated lower working scheme and makes no assertion of a new pair's
chain closure or chart clause.
\end{lemma}

\begin{proof}
Presentation arrows across levels are supplied by the construction of
the templates: the vertex term-hull LS witnesses give componentwise
strongness in the marked application, and the tile specialization of
\Cref{sys:tiling} gives the additional basis condition in the
geometric application.  Coherence is componentwise coherence, with the
finite independence test for geometric arrows.  Once a componentwise
union is an object, its arrows and its common-upper-bound property,
including into a larger object, follow from the AEC chain axioms and,
for geometry, finite character of independence.

Suppose an object $X$ at $\nu\in S$ maps to an object $Y$ at a larger
$\mu\in S$.  Present $Y$ on a well-order.  The image of $X$ uses at most
$\nu$ positions: each of its $\nu$ entries has finite term support, and
there are finitely many vertices.  Include these positions in a
well-ordered subpresentation of cardinality $\nu$, also including the
fixed named positions.  Coherence gives an arrow from $X$ to this
subpresentation.  This is mixed LS.

For a tuple bounded at $\nu$, the latter arrow preserves its orbit by
the same-level adequacy condition.  The chart clause identifies that
orbit with $D_F$.  Inclusion of the subpresentation in the large chart
leaves the labeled support pattern unchanged, so the same $D_F$ test
computes the diagram in $Y$.  This proves preservation along arbitrary
mixed arrows.  A bounded well-ordered pattern from a larger chart can
also be placed in an order of size $\nu$.  When its old tuple has the
same diagram as a specified tuple in a $\nu$-object, the same-level
chart clause gives an isomorphism matching those tuples.  Transporting
the enlargement through that isomorphism proves the required
mixed-level bounded saturation.

Root prescriptions pass along mixed arrows because the approximation
lists give the same majority comparisons there.  To compare the full
system types after projection, apply the type correspondence at the
receiving level to those equal root marginals.  Finally, the criterion
for maps obtained by exhausting bounded tuples is componentwise the
strong-submodel criterion of the working package.  In the geometric
case its additional condition is a finite independence comparison,
preserved by the matched system diagrams.  These arguments verify the
mixed clauses without inserting them as new proper-class parameters in
the unary adequacy formula.
\end{proof}

\subsection{A fixed register of bound requests}

We specify what is included when a closure requirement is invoked.
Given a uniform unary formula $S(\mu)$ and a bound request with ordinal
input $\alpha$ and descriptor $d$, write
\[
 B_q(d,\alpha;S)
 =\begin{cases}
 \text{the least ordinal satisfying the request},&
                         \text{if such an ordinal exists},\\
 0,&\text{otherwise}.
 \end{cases}
\]
In this notation $q$ belongs to the finite register below.  The notation
does not permit an arbitrary formula as a value of $q$, or an arbitrary
class as a parameter $S$.

\begin{lemma}[Uniform bound register]
\label{close:recipe-register}
All the refinements in one application of the one-dimensional tools,
one marked-increment step with its projection and boundary arguments,
or one geometric categoricity, basis, or boundary step are expressible
using a fixed finite register of first-order bound requests.  Their
parameters are finite descriptors, ordinal size inputs, and unary
adequacy tests from Proposition~\ref{close:uniform-definitions}.
Any references within one such step to earlier refinements can be
eliminated by a fixed finite number of substitutions, independently of
the input depth.  Every bound request needed on the levels where its
preceding hypotheses have been established has an ordinal witness.
\end{lemma}

\begin{proof}
Here is the register; its descriptions specify first-order tests, rather
than codes for formulas of increasing depth.

\paragraph{Cardinal, order-pattern, and partition bounds.}
Include successors, power sets, and their bounded iterations.  For
example, profile growth through $\alpha^+$ stages is bounded by the
recursion $a_0\ge\alpha$, $a_{\xi+1}=(2^{a_\xi})^+$, with suprema at
limits, for $\xi\le\alpha^+$.  This is a set recursion with ordinal
parameters.  Cut filling and saturation of order patterns of size at
most $\alpha$ use the same operations and the corresponding
$\alpha^+$-stage construction.  The finite-arity partition request is
the least ordinal on which every coloring of $k$-element subsets by the
given color set has a homogeneous subset of the prescribed size; $k$
and the size bounds are descriptor or ordinal inputs.  The partition
argument used to construct the order templates proves existence.
Small trees and their parameter accumulators use successors, powers,
and the least $\epsilon$ with $2^\epsilon>\chi$.  These operations cover
the chain lengths and cardinal costs in the domination and return-tree
arguments.

\paragraph{Placed comparison and syntax tests.}
For a specified permitted layout and tuple arity, the comparison request
asks for a bound on the size of a parameter test witnessing inequality
between two placed comparisons.  More explicitly, quantify over an
eligible level, its indicated order layout, the two candidate term lists
in order enlargements, and the parameters generated by the designated
parts.  If their comparisons differ on some bounded parameter list,
require that they differ on one of size at most the proposed bound.
Inequality is evaluated in a larger eligible padded presentation,
using isomorphism to a well-ordered chart, or equivalently by its $D_F$
comparison where the chart clause has been obtained.  All the objects
and positions involved in a particular test are sets.  Universal
quantification over the level cardinal is an ordinary first-order
quantifier.

The permitted layouts are the whole-model layout, designated vertex
parts, and finitely segmented or tiled layouts with the prescribed
order and colors.  Their codes belong to the descriptor set.  The
ordinary and rationally tiled ordinal-vector arguments of
\Cref{calc:testing,exact:tiled-testing} prove that this request has a
uniform witness.  The syntax request asks for a test size at which the fixed
$D_F$ comparison on patterns of a specified size is already determined.
The set of such patterns and the set-structure back-and-forth
construction in \Cref{lev:diagram-coding} prove existence.  Neither
request mentions a truth predicate for $\mathcal C_n$.

For the rationally tiled layouts $\sum_i(H_i\times\mathbb Q)$ used
in the geometric application, the comparison request includes the
bounded frozen data of Lemma~\ref{exact:tiled-testing}.
That lemma supplies its witness also for candidates in arbitrary
enlargements of the underlying $E$-order.  The exceptional fibers,
their frozen positions, and the set of cut coordinates have bounds
in the input arities and fixed descriptor; no new dependence on the
level or the finite depth is introduced.

\paragraph{Approximation, compactification, and isolation bounds.}
The approximation and compactification constructions take a comparison
bound and a profile bound at the specified arity, increase the lengths
to dominate the exception counts, and perform the bounded power-set
iteration just described.  Thus their sizes are fixed compositions of
the first two requests and cardinal operations.  The isolating-fragment
bound is the size of the initial supports, the bounded initial parts of
the candidate blocks, and the terms on those supports in the proof of
\Cref{iso:alternative}.  These again use the same comparison and
cardinal requests.  The construction of all-dominated objects uses
larger input ordinals in these requests, mixed LS, and the indicated
successor-length recursions; it introduces no new form of predicate
depending on the marked depth.  The existence assertions used here are
\Cref{iso:finite-existence,iso:all-lengths,iso:independent-lists}, with the mixed interface just verified.

\paragraph{Eventual orbit thresholds.}
The remaining request asks for the first threshold above which a
specified reversal, orbit-preservation, or stabilization assertion holds
at the eligible levels and the indicated arity.  The assertion itself
uses the standardized object sets, arrows, and orbit relation defined
above.  Thus the threshold test has the form
\[
 \forall\mu\bigl(S(\mu)\land\mu>\beta
           \ \Longrightarrow\ R(d,\alpha,\mu)\bigr),
\]
or its two-level version, with $R$ one of these fixed orbit assertions.
The reversal and stabilization arguments prove the existence of the
threshold when the current categoricity and the previous required
conditions hold on the unbounded eligible class.  The use for geometric
shapes has exactly those hypotheses by
\Cref{sys:working-package}.  Passing from these thresholds to
diagram homogeneity also requires smaller eligible levels cofinal
below the final level; that requirement is recorded separately as a
limit-point refinement.

We check that this list accounts for a whole generic step.  In a marked
step, extraction uses partition bounds, the underlying comparison and
isolation bounds, and arbitrarily large samples.  It chooses a template
from the set established in
Proposition~\ref{close:uniform-definitions}; see
\Cref{mark:template}.  After that choice, the two-part comparison of
\Cref{mark:two-part-testing}, used in \Cref{mark:chain-closure},
is the placed comparison request with
this template as parameter.  The downward isolation argument requires
the same bounds and sufficiently closed smaller levels.  Orbit
comparison and saturation use the transport argument; the passage to
the unary chart clause uses the syntax request.  Projection, completion,
and exact boundary lifting then use the ordinary comparison requests,
the prior unary conditions, and the tree costs; exact lifting may require
one further supply of cofinal smaller levels.  These are respectively
the constructions of \Cref{iter:projection-types,iter:boundary-completion,iter:exact-boundary}.

For geometry, first use the current categoricity condition and prior
basis conditions as unary inputs to the eventual orbit request, then
require smaller eligible levels and the chart syntax bounds.  Basis
lifting and exact boundary lifting use the same one-dimensional and
tiled-layout requests and tree costs, by
\Cref{sys:basis-classification,exact:geometric-boundary}.  A simultaneous projection calculation
or a maximum over finitely many lower descriptors is expressed by
bounded quantification over that finite descriptor.  It does not require
successively inserting their formulas into one another.

Within each of these generic proofs there are finitely many successive
uses of closure and of cofinal smaller levels.  Label those positions
in the proof in their order.  The predicate for a position is the initial
adequacy predicate conjoined with the earlier stated threshold,
closure, and limit-point conditions.  Substituting these descriptions
at the finitely many positions gives a fixed formula for each request.
The number of positions is attached to the generic proof, not to the
finite axis count; finite lower-scheme lists are handled by the bounded
quantification just described.  This produces the asserted finite
register with $S$ eliminated in favor of its uniform defining formula.

\paragraph{The eligible predicates in a marked step.}
The following substitution order also previews the common-level induction
at the end of this section.  References here to retaining $G_j$ use the
\hyperref[lev:closure-preview]{advance notation in Section~\ref*{lev:section-working-schemes}};
the actual classes and
their absorption property are constructed below, after the register is
fixed.  They supply levels for that induction, not parameters for the
requests being defined here.

Let $S_0$ be the
unary conjunction of the full individual working conditions for the
input scheme and the projection and stationary-prescription conditions
already established for the finitely many lower descriptors.  The mixed
interface supplies their mixed working conditions.  First refine $S_0$
by the testing, approximation, isolation, and domination bounds of
Sections~\ref{calc:section-calculus} and~\ref{iso:section-isolation};
call this class $S_1$.  On its sufficiently closed members the actual
constructions give the samples used in Lemma~\ref{mark:template}.
The induction maintains that $S_1$ is unbounded by retaining a $G_j$;
unboundedness of an arbitrary class alone would not justify this step.

After selecting $F^b$, let $R(\mu)$ be the conjunction of $S_1(\mu)$,
the individual presentation flag for $F^b$, and the natural-interpretation
flags for its base and upper values in the old scheme.  Evaluate the
testing request for the layout $U_H\cup A_{H'}$ on this preliminary
class $R$, using the old scheme's $D_F$.  Lemma~\ref{mark:template}
supplies the individual validity and natural interpretations;
Lemma~\ref{mark:two-part-testing} reduces the comparison to the old
working scheme, where the ordinal-vector argument of
Proposition~\ref{calc:testing} supplies its uniform bound.

Now refine $R$ by that testing bound and its placement and arity costs;
call the resulting class $S_2$.  Thus $S_2$ is the output of this
refinement, not the eligible predicate in its own defining request.
The preliminary predicate $R$ uses only the earlier costs and the two
specified kinds of flags.  Neither the request nor its witness uses
a $D_{F^b}$ chart clause or chain closure for the pair scheme.

At members of $S_2\cap\Lim(S_2)$ closed under the same arity costs,
Lemma~\ref{mark:chain-closure} proves pair-chain closure, including
bounded tuples spanning the chain.  The actual pre-closure orbit
argument and this chain assertion then give
Theorem~\ref{mark:package}; only now impose the syntax bound for
$D_{F^b}$ and request its chart clause.  Projection correspondence and
amalgamated boundary completion follow after the fixed comparison and
tree-cost refinements in Section~\ref{iter:section-boundaries}.
The fixed-parameter boundary argument and exact lifting may add one
further limit-point refinement to obtain smaller levels where
completion already holds.  Each displayed eligible predicate is thus
a finite conjunction of unary flags, previously listed bound closures,
and limit-point conditions.  Substitution follows this displayed order;
no predicate for a previously chosen $G_j$ or closure function is
inserted into a bound request.  More generally the auxiliary subclass
$T$ in the eligible-level convention is used only to ensure an
unbounded supply of levels through these finite refinements.  It is
not an additional parameter of any request.  In the present induction
$T=G_j$ supplies this existence assertion, while each request still
uses the unary adequacy predicate and the finitely expanded refinements
just displayed.

\paragraph{The eligible predicates in a geometric step.}
The first assertion of Theorem~\ref{sys:basis-classification} for the
current shape follows from the previously proved assertions for its
finitely many lower shapes.  Start with the unary conjunction of this
current categoricity assertion and the needed lower-shape assertions.
The structural and tiled-presentation proofs give the actual arrows,
coherence, union conditions, and mixed containment used by
Proposition~\ref{sys:working-package}.  Its finite-tile alignment and
orbit calculation prove eventual reversal on this class.  The first
route in Proposition~\ref{lev:scheme-diagram-routes} then proves
stabilization; impose those arity bounds and a limit-point refinement
within the current eligible class.  Actual diagram homogeneity is
proved at the retained members before the $D_F$ syntax bound is imposed.

The four subsequent claims in the proof of
Theorem~\ref{sys:basis-classification} use this established package,
the preceding one-dimensional bounds, and finite tree costs.  They
prove the relative-basis assertion at the retained levels.
For exact boundary lifting, include the already proved exact assertions
for all proper restricted shapes in the initial unary conjunction.
The tree and unoccupied-strip proofs of
Theorem~\ref{exact:geometric-boundary} add only the indicated layout,
support-closure, and tree-cost refinements.  Increasing the number of
axes changes the finite list of lower descriptors and the ordinal
size inputs; it does not add a new kind of eligible predicate or a new
position in this generic substitution order.

Finally, the witness assertions are precisely the uniform bounds proved
in the cited constructions with their stated hypotheses.  An irrelevant
descriptor, or a stage whose hypotheses fail, may give a request with
no witness.  Its value $0$ makes a total definable ordinal function and
has no mathematical force.  At a used stage we invoke the preceding
construction to establish existence before applying the bound.  In
particular, totalization supplies neither an adequate template nor a
missing lifting or testing witness.
\end{proof}

\subsection{The common majorant and its limit refinements}

We now construct the hierarchy used in the preceding previews.  The
register and its finite substitutions have been fixed independently of
this construction; the common majorant bounds their values, while the
limit refinements supply smaller eligible levels cofinally below a new
level.

Choose a cardinal $\vartheta$ above the hereditary sizes of the normalized
descriptor set and all fixed parameters.  This is one cutoff, chosen
before any finite-depth induction; it includes all candidate templates,
not just those already selected.  For $\alpha\in\Ord$, take the supremum
of the values of all the functions in the expanded finite register,
for every descriptor and for all ordinal inputs at most $\alpha$.
Finite lists of size inputs are included by taking all their coordinates
at most $\alpha$.  Each such collection is a set, and each request has
been totalized, so Replacement bounds its values.  We may therefore
choose a fixed, monotone definable ordinal function $\Phi$ such that
\[
 \Phi(\alpha)>\vartheta,\qquad \Phi(\alpha)>\alpha,
 \qquad B_q(d,\vec\beta)<\Phi(\alpha)
 \quad(\max\vec\beta\le\alpha).
\]
Include the initial single-model requirements in this finite register.
The formulas defining its members have already been fixed.  No partial
truth construction, and no bound on a variable formula's complexity,
is being used here.

For a definable class $A$ of ordinals write
\[
 \Lim(A)=\{\lambda:\sup(A\cap\lambda)=\lambda>0\},
 \qquad
 \operatorname{Cl}(f)=
 \{\lambda:\forall\alpha<\lambda\ f(\alpha)<\lambda\}.
\]
Multiple ordinal inputs can be replaced by their maximum in the closure
condition.  Define
\[
\begin{split}
G_0&=\{\mu:\mu\text{ is a strong limit cardinal},
                    \cf(\mu)=\omega,\ \mu\in\operatorname{Cl}(\Phi)\},\\
G_{j+1}&=G_j\cap\Lim(G_j)\qquad(j<\omega),\qquad
G_\omega=\bigcap_{j<\omega}G_j.
\end{split}
\]
Because $\Phi(0)>\vartheta$, every member of every $G_j$ exceeds the
common parameter cutoff.  This detail cannot be replaced by discarding
an arbitrary initial segment after each new template choice.

The hierarchy itself has a uniform definition.  For an ordinal $\lambda$,
form $A_0=G_0\cap(\lambda+1)$ by Separation and then, by finite recursion
on subsets of $\lambda+1$, put
\[
 A_{i+1}=\{\beta\in A_i:\beta>0\text{ and }
                                  \sup(A_i\cap\beta)=\beta\}.
\]
The assertion $\lambda\in A_j$ is one formula with $j,\lambda$ as inputs,
and is equivalent to $\lambda\in G_j$.  Quantifying over $j\in\omega$
defines $G_\omega$.  In particular the countable constructions below
can choose least witnesses using this formula and Replacement.

\begin{lemma}[Closure and absorption]
\label{close:absorption}
Each $G_j$ is unbounded and closed under suprema of strictly increasing
countable sequences of its members.  The class $G_\omega$ is unbounded.
If a current eligible class $S$ contains $G_j$, and a finite sequence of
refinements consists of the bound closures and thresholds in the fixed
register and restrictions to limit points of earlier eligible classes,
then the resulting class contains $G_{j+m}$ for some finite $m$.
Here every used bound must have a witness under its stated hypotheses.
\end{lemma}

\begin{proof}
Given any starting ordinal, choose strictly increasing infinite cardinals
$(\kappa_n)_{n<\omega}$ above it with
\[
 \kappa_{n+1}>
 \max\bigl(2^{\kappa_n},
                \sup\{\Phi(\alpha):\alpha<\kappa_n\}\bigr).
\]
The supremum $\kappa$ of this sequence is a cardinal of cofinality
$\omega$.  For every cardinal $\lambda<\kappa$, some $\kappa_n$ exceeds
$\lambda$, and then $2^\lambda\le2^{\kappa_n}<\kappa$.  For every
$\alpha<\kappa$, the same choice gives
$\Phi(\alpha)<\kappa_{n+1}<\kappa$.  Thus $\kappa\in G_0$.
If instead the $\kappa_n$ are an increasing sequence in $G_0$, their
supremum is again a countable-cofinality strong limit cardinal and is
strictly $\Phi$-closed by these same two tests.  This proves the two
assertions for $G_0$.

Suppose they hold for $G_j$.  Above any starting ordinal choose a strictly
increasing countable sequence in $G_j$.  Its supremum belongs to $G_j$
by closure and to $\Lim(G_j)$ because the sequence is cofinal there.
This proves unboundedness of $G_{j+1}$.  A countable increasing sequence
in $G_{j+1}$ has its supremum in $G_j$, and its own members witness that
this supremum belongs to $\Lim(G_j)$.  Hence $G_{j+1}$ has the asserted
closure as well.

Choose $\lambda_n\in G_n$, strictly increasing and above an arbitrary
prescribed ordinal.  For each fixed $j$, its tail lies in $G_j$ because
the classes are nested.  The supremum belongs to $G_j$ by its countable
closure.  It consequently belongs to $G_\omega$, proving unboundedness.

For absorption, maintain that the current eligible class contains a
particular $G_k$.  A request value on inputs below a member $\mu\in G_k$
is less than $\Phi(\alpha)$ for a suitable $\alpha<\mu$, and hence less
than $\mu$.  Thus a bound-closure refinement retains all of $G_k$.
A constant threshold is included as a request with input $0$ and has
the same property.  If the next refinement demands cofinally many
members of the current class below the new level, retain $G_{k+1}$:
its members are limit points of $G_k$ and thus of the current class.
If several earlier eligible classes are mentioned, the maintained
index can first be increased to the maximum of their finitely many
indices.  Induction through the finite refinement list proves the
claim, with $m$ bounded by its number of limit-point refinements.
\end{proof}

\begin{theorem}[Common levels for all finite depths]
\label{close:common-levels}
With the fixed data of either application and the constructions proved
in the preceding sections, there is an unbounded class of
countable-cofinality strong limit cardinals at which all its finite-depth
conclusions hold simultaneously.  In the marked application this
includes the working packages, projection and type correspondences,
amalgamated completion, and exact boundary lifting for every
$\mathcal C_n$, uniformly over the root prescription codes realized
at arbitrarily large eligible single-model levels.
In the geometric application it includes categoricity, the working
packages, basis lifting, and exact boundary lifting for every finite
shape.  One such common level may be chosen above any specified set of
fixed cardinal thresholds.
\end{theorem}

\begin{proof}
Use the fixed register and $\Phi$ just constructed.  The initial
single-model results hold on $G_0$ after the initial requirements have
been included.  Suppose the assertions required as input to a finite
induction step hold throughout $G_j$.  Take $S$ to be the class defined
by their unary adequacy formula, with the current descriptor.  It
contains $G_j$.  Lemma~\ref{close:mixed-interface} supplies the mixed
conditions for that step.  Lemma~\ref{close:recipe-register} identifies
its bound requests and gives the existence of each relevant witness
at its place in the proof.  Lemma~\ref{close:absorption} then retains
some $G_{j'}$, with $j'<\omega$, through all its refinements.

When a marked step extracts a new template, its code lies in the fixed
descriptor set by Proposition~\ref{close:uniform-definitions}.  The
remaining requests with that code as parameter were already majorized
by $\Phi$.  Thus making the choice introduces no new closure function
or cutoff.  For a given finite shape, the finitely many previously
required lower shapes can first be brought to one $G_j$ by taking the
largest of their indices.  The same observation applies to the finite
projection lists in the marked argument.  The two ordinary finite
inductions therefore give, for each finite descriptor, all its required
conclusions on $G_j$ for some finite $j$.

In the marked application only root codes consistent at arbitrarily
large eligible single-model levels are used.  A code obtained from an
actual pair at a sufficiently closed starting cardinal has that property:
freely lift its type with the named approximation data to arbitrarily
large single-model levels.  The majority rule is preserved there, and
restriction from such a level gives the same stationary prescription
at every sufficiently closed smaller single-model level with that
named pattern.  This is the prescription argument in
Lemma~\ref{geom:root-prescription}, prior to the use of the finite marked
systems.  It requires no persistence of a no-growth statement about
$P$ at larger levels.  The fixed code set and majorant include every
such possible input code, so the common-level construction can precede
the choice of the actual pair.

Every member of $G_\omega$ belongs to $G_j$ for every finite $j$.  It consequently
satisfies the conclusion for every finite descriptor used in the
application, although the index $j$ and the selected template may depend
on that descriptor.  Lemma~\ref{close:absorption} makes $G_\omega$
unbounded, so choose a member above the desired thresholds.  All
constructions in this argument are ordinary set recursions, Separation,
Replacement, and finite induction with uniform formulas.  This completes
the simultaneous closure justification in ZFC.
\end{proof}
